\documentclass{article}
\usepackage{graphicx} % Required for inserting images

\title{Sutton and Barton Exercises Chapter 8}
\author{danieltocco0 }
\date{December 2025}

\begin{document}

\maketitle

\section{Chapter 8}
\textit{Exercise 8.1} The nonplanning method looks particularly poor in Figure 8.3 because it is
a one-step method; a method using multi-step bootstrapping would do better. Do you
think one of the multi-step bootstrapping methods from Chapter 7 could do as well as
the Dyna method? Explain why or why not.
\\
\\

In short, no. Only for states that are visited can be updated. To take the 

extreme of Monte Carlo, for example, only n visited states would be updated. 

1700 - n states would still need to be updated thereafter.
\\\\
\textit{Exercise 8.2} Why did the Dyna agent with exploration bonus, Dyna-Q+, perform
better in the first phase as well as in the second phase of the blocking and shortcut
experiments? 
\\
\\

Due to the r + $\kappa\sqrt{\tau}$ term, Dyna-Q+ is more explorative than Dyna. 

It is expected Dyna-Q+ would thus find the optimal path faster for both 

phases. Similarily to phase two, Dyna in phase one may possibly even find 

a sub-optimal path from which it cannot escape.
\\\\
\textit{Exercise 8.3} Careful inspection of Figure 8.5 reveals that the difference between Dyna-Q+
and Dyna-Q narrowed slightly over the first part of the experiment. What is the reason
for this?
\\
\\

Because Dyna and Dyna-Q+ both converge to an optimal solution, but 

Dyna-Q+ retains it's stronger exploration. To elaborate when Dyna-Q+ 

explores a bad (relative to phase) path, Dyna-Q is still behaving optimally. 

This closes the gap until Dyna-Q goes back to it's optimal path. 
\\\\
\textit{Exercise 8.5} How might the tabular Dyna-Q algorithm shown on page 164 be modified
to handle stochastic environments? How might this modification perform poorly on
changing environments such as considered in this section? How could the algorithm be
modified to handle stochastic environments and changing environments? 
\\
\\

In step e), use the expected reward based on the samples we have. The 

problem here is that the environment could change, thus changing our 

expectation. The expectation may very slowly update, however. What we

want to do is add a recency weighing-term like: $P_{new} = (1 - \beta) * P_{old} + \beta *$ 

indicator$(new_{transition})$ for the probabilities, or something like Dyna-Q+,

where, again, we encourage the agent to try long untested actions which 

would discover changes in the environment.
\\\\
\textit{Exercise 8.6} The analysis above assumed that all of the b possible next states were
equally likely to occur. Suppose instead that the distribution was highly skewed, that
some of the b states were much more likely to occur than most. Would this strengthen or
weaken the case for sample updates over expected updates? Support your answer.
\\
\\

It would strengthen the case for sample updates. Regardless of the

distribution, expected updates must account for ALL transitions. 

A skewed distribution for sample updates would reduce the variance, 

thus making it even better.
\\\\
\textit{Exercise 8.7} Some of the graphs in Figure 8.8 seem to be scalloped in their early portions,
particularly the upper graph for b = 1 and the uniform distribution. Why do you think
this is? What aspects of the data shown support your hypothesis? 
\\
\\

Because for b = 1, you can only possibly transition to one state. The starting 

values are far from their actual values, and combine this with the 

stochasticity of action-selection (even if there are only two, but

reward dynamics must be taken into account), you can undershoot or 

overshoot by a good amount at the start. Over many episodes, variance 

will decrease, hence the curve smoothens. Following similar logic, 

larger branching factors will have less variance, hence they don't 

scallop to the same extent.
\end{document}
